\section{Conclusion and Future Work}
\label{sec:conclusion}

In this paper, we conducted a rigorous, end-to-end investigation into the orthographic robustness of Roman Urdu sentiment classification under authentic, human-documented spelling variations. By performing an exhaustive data audit on the 28,090-review RUWV-NSR dataset, we canonicalized eight corrupted label strings into five classes and eliminated 4,632 exact duplicate reviews (16.49\%), ensuring a strictly leakage-free benchmark split.

We evaluated classical (TF-IDF + Logistic Regression) and deep multilingual architectures (mBERT, XLM-RoBERTa full fine-tuning, and XLM-RoBERTa LoRA adaptation) across canonical noise tiers (Clean, Mild, Moderate, Severe) and a $3 \times 3$ factorial grid disentangling substitution rate from edit distance. Our empirical findings demonstrate:
\begin{enumerate}
    \item \textbf{LoRA Superiority in Robustness:} Parameter-efficient adaptation of XLM-RoBERTa via LoRA demonstrated remarkable resilience, suffering only a 4.68\% relative Macro-F1 drop under severe noise compared to 7.01\% for full fine-tuning and 16.42\% for TF-IDF.
    \item \textbf{Factorial Dynamics:} Word substitution density accounts for over 82\% of performance variance, demonstrating that the sheer frequency of corrupted tokens dominates character edit distance in determining classification failure.
    \item \textbf{Subword Fragmentation Impact:} Although XLM-RoBERTa's larger SentencePiece vocabulary yields fewer absolute subwords per token (2.28 vs.\ 2.47 for mBERT), the proportional fragmentation increase under perturbation is comparable (27.5\% vs.\ 24.3\% of pairs affected). This confirms that LoRA's robustness advantage stems primarily from its implicit regularization rather than tokenization geometry.
    \item \textbf{Mitigation Pathways:} Generic text preprocessing yields zero or negative recovery, whereas lexicon-driven normalization recovers 84.23\%--93.79\% of lost performance. Furthermore, 50\% noise-augmented training with LoRA elevates clean and noisy Macro-F1 simultaneously.
    \item \textbf{Cross-Domain Generalizability:} Findings replicated consistently across 3,000 independent reviews from the Daraz.pk e-commerce corpus.
\end{enumerate}

In future work, we plan to extend this evaluation framework to phonetic embedding alignments, explore phoneme-aware contrastive pre-training, and benchmark emerging large language model (LLM) instruction tuning across South Asian code-mixed dialects.
